\section{L1 Introduction \& Task Decomposition}

\begin{definition}{Vocabulary}\\
    Explain Host, Device, CU, PE, FU, UE, task.
    \tcblower
    Host delegates work to a Compute \important{Device}; device = CU(s) $\to$ PE(s) $\to$ FU(s). 
    UE = process/thread that runs a task; task = OS process/thread \emph{or} a piece of parallel work.
    %TODO: insert image slide 12 with notes
\end{definition}

\begin{example2}{Ex.\ S13 -- CB1 / Jetson terminology}\\
    Programming on the CB1, delegating to the Jetson Nano -- what is what?
    \tcblower
    CB1 = \important{Host}; Nano = \important{Device}; Nano SMs = CUs; CUDA cores = PEs; ALU/FPU = FU; spawned threads = UEs; one Sobel slice = task.
    %TODO: insert translation from markdown
\end{example2}

\begin{KR}{Recipe: Map Platform $\to$ Terminology}\\
    Host (delegates) $\to$ Device (receives) $\to$ CU $\to$ PE $\to$ FU. UE = process/thread running a task.
\end{KR}

\subsection{Sobel}

\begin{code}{Sobel Algorithm}\\
    %TODO: insert image from annotated slides!!
\end{code}

\begin{formula}{Sobel Operator}\\
    $G_x=\begin{bsmallmatrix}-1&0&+1\\-2&0&+2\\-1&0&+1\end{bsmallmatrix}$,\;
    $G_y=\begin{bsmallmatrix}-1&-2&-1\\0&0&0\\+1&+2&+1\end{bsmallmatrix}$;\;
    $|G|=\sqrt{G_x^2+G_y^2}\to|G_x|+|G_y|$
\end{formula}

\begin{example2}{How does Sobel work?}\\
    Determine the edge-detection computation.
    \tcblower
    Two $3\times3$ convolutions $G_x,G_y$ give the gradient; magnitude $|G|=\sqrt{G_x^2+G_y^2}$ \important{simplified} to $|G_x|+|G_y|$, then threshold.
\end{example2}

\subsection{Parallel P}

\begin{KR}{Recipe: Task-Decompose an Algorithm}\\
    \textbf{Draw the CFG} nodes = functions/actions; directed edges = data dependency (producer$\to$consumer).

    \textbf{Count distinct actions} = number of functional tasks.

    \textbf{Find independent tasks} no path between them $\Rightarrow$ may run concurrently.

    \textbf{Try both splittings} functional (by named function, few) vs loop/data (by region, scales -- prefer).

    \textbf{State parallelism + joins} how many overlap, where they synchronise.
\end{KR}

\begin{KR}{Recipe: Dependency Analysis}\\
    \textbf{Group} tasks sharing data/constraints.

    \textbf{Order} add constraint where one consumes another's output (+ ghost-exchange before boundary).
    
    \textbf{Classify shared data} read-only / effectively-local / read-write (accumulate, multi-read-single-write).
\end{KR}



\begin{KR}{Recipe: Evaluate a Parallel Design}\\
    Check: compute $>$ management overhead; \#tasks $\geq$ \#PEs; data dependencies minimal; schedule (static if equal, else dynamic / work-stealing).
\end{KR}







\begin{example2}{Ex.\ S22/23/25 -- Pattern language}\\
    Name the pattern hierarchy and the Finding-Concurrency patterns.
    \tcblower
    Hierarchy: Finding Concurrency $\to$ Algorithm Structure $\to$ Supporting Structure $\to$ Implementation Mechanisms. Finding Concurrency = \important{Decomposition} (task/data) + \important{Dependency Analysis} (group/order/share).
\end{example2}

\begin{example2}{Ex.\ S26/27 -- Decompose Sobel}\\
    Split Sobel into tasks; identify parallelism; functional or loop splitting?
    \tcblower
    6 functional tasks; \texttt{sobel\_x}$\parallel$\texttt{sobel\_y} independent. \important{Both} apply: functional (limited, 2-way), loop splitting (region-split, scalable -- preferred).
\end{example2}

\begin{example2}{Ex.\ S28/29 -- Read the CFG}\\
    Distinct actions? Independent tasks? Parallelisation potential? What is a CFG?
    \tcblower
    \important{6} distinct; only \texttt{sobel\_x},\texttt{sobel\_y} independent; functional parallelism limited (2-way), data decomposition scales. CFG = control-flow graph: all execution paths; used by compilers/profilers and to spot concurrency.
    
    \includegraphics[width=\linewidth]{sobel_cfg.png}
\end{example2}

\begin{example2}{Ex.\ S30--33 -- Dependency analysis}\\
    Which sub-pattern is most relevant? Group/order/share for Sobel?
    \tcblower
    \important{Data Sharing} is the interesting one. Group: \texttt{sobel\_x/y} (same grey input). Order: \texttt{sobel\_all} after both finish. Share: grey = read-only, kernel = effectively-local, \texttt{Gx/Gy} = read-write, \texttt{sobel\_all} = \important{accumulate}.
\end{example2}

\begin{example2}{Ex.\ S34--37 -- Evaluate task parallelism}\\
    Compute vs management? Tasks vs PEs? Data deps? Schedule/load balancing?
    \tcblower
    Whole-function task justifies a UE (per-pixel doesn't). Want \#tasks $\geq$ \#PEs $\to$ region-split. Only dep is the \texttt{sobel\_all} join. Equal regions $\to$ \important{static}; else dynamic / work-stealing.
\end{example2}

\begin{example2}{Ex.\ S38/39 -- Fork--Join}\\
    What is fork-join? When is it suitable?
    \tcblower
    Clone code into concurrent tasks, then \emph{join} (wait). \important{Not} POSIX \texttt{fork()}. Suits recursive/irregular/dynamic task sets; a poor fit for a fixed regular pipeline like Sobel.
\end{example2}

\begin{example2}{Ex.\ S41--45 -- Process vs Thread}\\
    What is a process? How do processes/threads differ?
    \tcblower
    \emph{Process}: running program; resource + scheduling unit; own \texttt{main()}/address space; MPU-isolated; \texttt{fork()} = CoW clone, parent \texttt{wait()}s. \emph{Thread}: inside a process, shares memory, cheap $\to$ reuse via \important{thread pool}.
\end{example2}
